\section{Tiered access and evidence reporting}
\label{sec:retrieval}
\label{sec:lanes}

A query can reach the evidence through two entry points. Catalogue-first
search selects a collection, examines shelf descriptions and opens selected
documents or passages. Pinned search starts with a supplied collection
identifier and returns candidate chunks through the selected retrieval path.
Both paths retain source identifiers for locating returned candidates.
Figure~\ref{fig:workflow} summarises the broader acquisition-to-report path;
this section focuses on its retrieval and evidence-reporting stages.

The workflow separates obtaining candidate evidence from reporting the outcome.

\begin{figure}[t]
\centering
\includegraphics[width=\columnwidth]{fig_tiered_workflow_vector}
\caption{Tiered access workflow. Pinned search and catalogue navigation
supply candidate references. The workflow can read candidates directly or
use Laya to admit passages before checking source text. When an interruption
is detected and the workflow reaches its reporting step, the agent is
instructed to report incomplete status, searched scope and available evidence;
an endpoint timeout may terminate the turn before that report is produced.}
\label{fig:tiered-algorithm}
\end{figure}

The chat endpoint sets timeouts and turn limits. Workflow instructions
specify the retrieval sequence and batch sizes and ask the agent to report
an exhausted call budget. An endpoint-level timeout can terminate the turn
before the agent produces a final scope or evidence report.

\subsection{Tier policies and the evidence gate}
\label{sec:trace}
\label{sec:gate}
\label{sec:gate-primitive}
\label{sec:twostage}

The knowledge-base chat interface defines two policies. The search policy
can pass a collection identifier to two-stage or vector search; a call
without an effective identifier may search across collections
(Section~\ref{sec:exp-e2}). Hierarchical resolution and optional graph
expansion also affect which candidates are returned, so the requested scope
is not a server-enforced access boundary. The librarian policy instead
starts from the full catalogue, uses shelf descriptions to choose document
references and directs the agent to score them in sequential batches before
reading the survivors. The two-stage tool uses lexical search, with optional
graph expansion, to form a candidate pool; dense ranking then returns
candidates for inspection.

Separately, the MCP server provides \path{kb_hybrid_search}, which combines
vector and catalogue candidates. When enabled, its gate sorts retained
documents by admission score (the verifier score returned by Laya), as in the
verifier-score ordering replay. The replay pools were built from up to 30 vector candidates per collection at a
cosine threshold of 0.35, with catalogue candidates added and the fused pool
capped at twelve (Section~\ref{sec:e6}).

The local Laya classifier \cite{laya} is exposed as
\path{kb_laya_judge}. In reference mode, it fetches candidate text, segments
it and assigns each passage an admission score in $[0,1]$ for the evidence
question. Passages meeting the threshold (0.5 by default) are eligible for
the returned pack. The gate orders eligible passages by source and position, removes duplicates
and applies its character limit. It returns coordinates for admitted source
spans alongside separate fetch and scoring errors, so valid passages can
remain available after a partial failure. A search-result chunk locator, a
gate-reported source span and the text ultimately included in a capped prompt
refer to different stages; a locator alone does not show how much evidence
reached the agent. This score governs passage admission; the workflow or
hybrid tool controls candidate fusion and document order. Figure~\ref{fig:tiered-algorithm} and
Algorithm~\ref{alg:tiered} distinguish direct reading from scoring, and an
incomplete review from a completed review that found no support.

\begin{table*}[t]
\centering
\small
\setlength{\tabcolsep}{5pt}
\caption{Access policies over the tiered path. Modes differ in how candidate collections and documents are selected.
The hybrid row describes the experimental composition used in the evaluation.}
\label{tab:lanes}
\begin{tabular}{@{}>{\raggedright\arraybackslash}p{0.11\textwidth}>{\raggedright\arraybackslash}p{0.19\textwidth}>{\raggedright\arraybackslash}p{0.24\textwidth}>{\raggedright\arraybackslash}p{0.16\textwidth}>{\raggedright\arraybackslash}p{0.20\textwidth}@{}}
\toprule
Policy & Tiers traversed & Retrieval affordance & Budget type & Intended terminal report\\
\midrule
Search & pinned collections, segment & Two-stage or vector search & Endpoint turn limit & Agent answer or reported scoped not-found\\
Librarian & catalogue, shelf, document, segment & Catalogue navigation and passage scoring & Workflow instructions & Agent answer or reported searched scope\\
Hybrid & catalogue and segment & Combined vector and catalogue candidates & Configuration-specific candidate cap & Candidates or judged survivors\\
\bottomrule
\end{tabular}
\end{table*}

Response transport also determines which candidates the agent can see.
A recorded transport incident truncated a document listing before the
agent could inspect the whole shelf. The shelf tool now caps its response and
supplies grouped and paginated views, allowing the caller to request the
remaining entries. Section~\ref{sec:eval} reports the repair check.

In this paper, evidence tracing means retaining document, part and span
coordinates for retrieved text, together with relevant fetch or scoring
errors. These references locate text for inspection; they are not immutable,
versioned citations and do not establish claim-level entailment. Durable
resolution after rename, reparse or replacement would additionally require a
content fingerprint or version identifier to be carried in the trace; that
property is not evaluated here.

\begin{algorithm}[tbp]
\caption{Tiered access with optional evidence gate}
\label{alg:tiered}
\begin{algorithmic}[1]
\Require question $q$; catalogue $\mathcal{C}$; threshold $\tau$; gate flag $g$
\State $D \gets \textsc{GatherCandidates}(q,\mathcal{C})$ \Comment{chosen policy}
\If{$g$}
  \State $S \gets \textsc{LayaJudge}(D,q,\tau)$ \Comment{admit passages}
  \State Optionally check selected source spans and retain their references
\Else
  \State $S \gets \textsc{ReadSelected}(D)$
\EndIf
\If{execution incomplete and the reporting step is reached}
  \State \Return incomplete status, scope and available evidence
\ElsIf{execution incomplete}
  \State \Return interruption status
\ElsIf{completed review finds no support}
  \State \Return scoped not-found with blind spots
\Else
  \State \Return answer with source references
\EndIf
\end{algorithmic}
\end{algorithm}

\subsection{Reporting search outcomes}
The workflow instructs the agent to cite sources when the evidence supports
an answer. After the agent reaches the reporting step for a completed review
without support, it reports the searched scope and remaining blind spots. If a
failure or timeout prevents that step, the outer endpoint may return an
interruption status without scope or evidence. Section~\ref{sec:eval}
evaluates these outcomes in observed agent trials.
